\documentclass[11pt]{article}

\usepackage[margin=1in]{geometry}
\usepackage{amsmath,amssymb,amsthm,amsfonts}
\usepackage{mathtools}
\usepackage{enumitem}
\usepackage{bm}

\newtheorem{theorem}{Theorem}[section]
\newtheorem{lemma}[theorem]{Lemma}
\newtheorem{corollary}[theorem]{Corollary}
\newtheorem{proposition}[theorem]{Proposition}

\theoremstyle{definition}
\newtheorem{definition}[theorem]{Definition}
\newtheorem{remark}[theorem]{Remark}
\newtheorem{algorithm}[theorem]{Algorithm}

\DeclareMathOperator{\Lop}{\mathcal{L}}
\DeclareMathOperator{\ord}{ord}
\DeclareMathOperator{\Res}{Res}

\title{A Universal Solution Methodology for Constant-Coefficient
Fractional Riccati Equations via the Polynomial Chebyshev--Wheeler
Recurrence}
\author{}
\date{}

\begin{document}
\maketitle

\section{Setup}

Fix
\[
P(u),\, Q(u),\, R(u) \in \mathbb{C}[u], \qquad \mu \in (0,1],
\]
and consider
\[
D^{\mu} y(t,u) = P(u) + Q(u)\, y(t,u) + R(u)\, y(t,u)^{2},
\qquad I^{1-\mu} y(0,u) = 0.
\]
All constructions are over the coefficient ring $\mathbb{C}[u]$.

\section{M\"untz--Tau coefficient sequence}

\begin{lemma}[M\"untz--Tau expansion]\label{lem:muntz}
The unique solution admits
$y(t,u) = \sum_{k\geq 1} a(k,u)\, t^{k\mu}$ with
$a(k,u) \in \mathbb{C}[u]$ given by
\[
a(1,u) = \frac{P(u)}{\Gamma(\mu+1)},
\qquad
a(k+1,u) = \frac{\Gamma(k\mu+1)}{\Gamma((k+1)\mu+1)}
\Bigl( Q(u)\, a(k,u) + R(u) \!\!\sum_{\substack{j+\ell=k\\ j,\ell\geq 1}}\!\!
a(j,u)\, a(\ell,u) \Bigr).
\]
\end{lemma}

\begin{proof}
Substitute the ansatz into the Volterra form
$y = I^{\mu}[P + Qy + Ry^{2}]$ and use
$I^{\mu} t^{s} = \Gamma(s+1)/\Gamma(s+\mu+1)\, t^{s+\mu}$. Match
$t^{(k+1)\mu}$ on both sides.
\end{proof}

\begin{lemma}[Polynomial degree growth]\label{lem:degree}
$\deg_{u} a(k,u) \leq k\, \max\{\deg P, \deg Q, \deg R\}$.
\end{lemma}

\begin{proof}
Induction on $k$ using the recurrence of Lemma~\ref{lem:muntz}. The
base case $k=1$ is $\deg_{u} a(1,u) = \deg P$. The inductive step is
controlled by the larger of $\deg Q + \deg a(k,u)$ and
$\deg R + 2\,\max_{j<k} \deg a(j,u)$.
\end{proof}

\begin{lemma}[Convolution closure]\label{lem:conv}
The bilinear convolution
$(a \ast a)(k,u) := \sum_{j+\ell=k,\, j,\ell\geq 1} a(j,u)\, a(\ell,u)$
preserves $\mathbb{C}[u]$ coefficient-wise. The generating series
$g(z,u) := \sum_{k\geq 1} a(k,u)\, z^{k}$ satisfies the algebraic
identity
\[
g(z,u) = P(u)\,\Gamma_{\mu}(z) + Q(u)\,(\Gamma_{\mu}\!\ast\! g)(z,u)
+ R(u)\,(\Gamma_{\mu}\!\ast\! g \ast g)(z,u),
\]
where $\Gamma_{\mu}$ encodes the Gamma-ratio weights.
\end{lemma}

\begin{proof}
Direct: the convolution is a finite $\mathbb{C}[u]$-bilinear sum. The
generating identity is the $z$-transform of Lemma~\ref{lem:muntz}.
\end{proof}

\section{Moment functional}

\begin{definition}[Moment functional]\label{def:Lop}
$m(k,u) := a(k+1,u)$ for $k\geq 0$;
$\Lop : \mathbb{C}[x] \to \mathbb{C}[u]$ by $\Lop[x^{k}] = m(k,u)$,
extended $\mathbb{C}[u]$-linearly.
\end{definition}

\begin{definition}[Hankel determinants]\label{def:Hankel}
$H_{n}(u) := \det\bigl[ m(i+j,u) \bigr]_{0\leq i,j \leq n-1} \in \mathbb{C}[u]$.
\end{definition}

\begin{lemma}[Quasi-definiteness]\label{lem:quasidef}
$H_{n}(u) \neq 0$ as an element of $\mathbb{C}[u]$ for all $n\geq 1$
under the genericity hypothesis $P(u)\, R(u) \not\equiv 0$. The set
$\{u \in \mathbb{C} : H_{n}(u) = 0\text{ for some }n\}$ is a countable
union of proper algebraic subvarieties of $\mathbb{C}$.
\end{lemma}

\begin{proof}
$m(0,u) = a(1,u) = P(u)/\Gamma(\mu+1) \not\equiv 0$, so $H_{1}\neq 0$.
For $n\geq 2$, the leading coefficient of $H_{n}(u)$ as a polynomial in
$u$ is computed from the highest-degree contributions to each $a(k,u)$,
which are themselves nonzero polynomials in $\deg P, \deg Q, \deg R$
under genericity. The vanishing locus of each $H_{n}$ is an algebraic
hypersurface in $\mathbb{C}$, hence finite.
\end{proof}

\begin{lemma}[Signed --- never positive --- definiteness]\label{lem:signed}
For every real $u$ at which the Volterra contraction hypotheses hold,
$\Lop$ is not positive-definite: there exists $p \in \mathbb{R}[x]$
with $\Lop[p^{2}] \notin \mathbb{R}_{\geq 0}$.
\end{lemma}

\begin{proof}
By Lemma~\ref{lem:conv} and the contraction hypothesis, $g(\cdot,u)$
is meromorphic with all poles strictly off $\mathbb{R}$. If $\Lop$
were positive-definite, Favard's theorem applied to the Jacobi matrix
$J$ of \S\ref{sec:jacobi} would force $\sigma(J) \subset \mathbb{R}$;
its companion-matrix realization gives the zeros of $P(M,\cdot)$ as
real, hence the poles of $\hat R(M,\cdot,u)$ as real, contradicting
the disjointness $\Delta_{g}(u) \cap \mathbb{R} = \emptyset$.
\end{proof}

\section{OPS and its generating polynomials}

\begin{definition}[Monic OPS]\label{def:OPS}
$P(n,x) \in \mathbb{C}[u][x]$, monic of degree $n$, satisfies
$\Lop[P(n,x)\, x^{k}] = 0$ for $0\leq k < n$, with
\[
P(n+1,x) = (x - \alpha(n))\, P(n,x) - \beta(n)\, P(n-1,x),
\quad P(-1,x)=0,\; P(0,x)=1.
\]
\end{definition}

\begin{definition}[Generating polynomial]\label{def:Sn}
$S(n,z) := \sum_{k\geq 0} \Lop[\, x^{k} P(n,x)\,]\, z^{k}
\in \mathbb{C}[u][[z]]$.
\end{definition}

\begin{lemma}[Vanishing order]\label{lem:vanord}
For $n\geq 0$,
$\ord_{z} S(n,z) = n$ and $[z^{n}] S(n,z) = h(n) := \Lop[P(n,x)^{2}]
\in \mathbb{C}[u]^{\times}$ on the complement of
$\bigcup_{k\leq n}\{H_{k}=0\}$.
\end{lemma}

\begin{proof}
Orthogonality forces $\Lop[x^{k} P(n,x)] = 0$ for $k < n$. At $k=n$,
$\Lop[x^{n} P(n,x)] = \Lop[P(n,x)^{2}]$ since $P(n,x)$ is monic and
lower-order terms are killed by orthogonality.
\end{proof}

\begin{lemma}[Christoffel--Darboux norm ratio]\label{lem:hratio}
$h(n) = H_{n+1}(u)/H_{n}(u)$, with $H_{0} := 1$. Hence
$h(n) \in \mathbb{C}(u)$, and $h(n) \in \mathbb{C}[u]^{\times}$
generically by Lemma~\ref{lem:quasidef}.
\end{lemma}

\begin{proof}
Standard: expand $P(n,x) = \det(\,\cdot\,)/H_{n}$ via Heine's formula
and compute $\Lop[P(n,x)^{2}]$.
\end{proof}

\begin{lemma}[Base $S(0,z)$]\label{lem:S0}
$S(0,z) = \sum_{k\geq 0} m(k,u)\, z^{k}$, with $S(0,0) = m(0,u) = P(u)/\Gamma(\mu+1)$.
\end{lemma}

\section{Polynomial Chebyshev--Wheeler recurrence}

\begin{theorem}[Polynomial Chebyshev--Wheeler]\label{thm:CW}
The sequence $\{S(n,z)\}_{n\geq -1}$ obeys
\begin{equation}\label{eq:CW}
S(n+1,z) \;=\; \frac{S(n,z) - S(n,0)}{z} \cdot \mathbf{1}_{n=0}
\;+\; \frac{S(n,z)}{z}\cdot \mathbf{1}_{n\geq 1}
\;-\; \alpha(n)\, S(n,z) \;-\; \beta(n)\, S(n-1,z),
\end{equation}
with $S(-1,z) := 1$ and unique $\alpha(n),\beta(n) \in \mathbb{C}[u]$
fixed by the divisibility constraint $\ord_{z} S(n+1,z) \geq n+1$.
Equivalently, in $\mathbb{C}[u][[z]]$,
\[
\beta(n) \;=\; \frac{[z^{n-1}](S(n,z)/z)}{[z^{n-1}] S(n-1,z)}
\;=\; \frac{h(n)}{h(n-1)},
\qquad
\alpha(n) \;=\; \frac{[z^{n}](S(n,z)/z) - \beta(n)\,[z^{n}] S(n-1,z)}
{[z^{n}] S(n,z)}.
\]
\end{theorem}

\begin{proof}
The three-term recurrence on $P(n,x)$ gives
$x P(n,x) = P(n+1,x) + \alpha(n) P(n,x) + \beta(n) P(n-1,x)$.
Apply $\Lop[\,\cdot\, x^{k}]$ and multiply by $z^{k}$:
$\Lop[x^{k+1} P(n,x)] = \Lop[x^{k} P(n+1,x)]
+ \alpha(n) \Lop[x^{k} P(n,x)] + \beta(n) \Lop[x^{k} P(n-1,x)]$.
Summing over $k\geq 0$ produces \eqref{eq:CW} after reindexing
$\sum_{k\geq 0} \Lop[x^{k+1}P(n,x)] z^{k} = z^{-1}(S(n,z) - S(n,0))$.
For $n\geq 1$, $S(n,0)=0$ by Lemma~\ref{lem:vanord}, so the shift is
genuine. The constraint $\ord_{z} S(n+1,z) \geq n+1$ has exactly two
linear equations over $\mathbb{C}[u]$ (matching $z^{n-1}$ and $z^{n}$),
solved as displayed.
\end{proof}

\begin{corollary}[Constructive recurrence in $\mathbb{C}[u][z]$]\label{cor:constructive}
At step $n$, both $S(n,z)$ and $S(n-1,z)$ are known polynomials
truncated to degree $\leq 2N$ for chosen working bandwidth $N$.
$\beta(n)$ is the $z^{n-1}$-coefficient match, $\alpha(n)$ the
$z^{n}$-coefficient match. $S(n+1,z)$ assembles by polynomial linear
combination. Total cost per step: $O(N)$ polynomial coefficient
operations in $\mathbb{C}[u]$. No Hankel matrix is formed; no power-series
inversion is performed.
\end{corollary}

\begin{corollary}[Standard-form coefficients]\label{cor:ABC}
In the form $P(n+1,x) = (A(n)x + B(n)) P(n,x) - C(n) P(n-1,x)$:
$A(n)=1$, $B(n) = -\alpha(n)$, $C(n) = \beta(n) = h(n)/h(n-1)$.
\end{corollary}

\begin{theorem}[Equivalence to Hankel form]\label{thm:hankelequiv}
On the complement of the discriminant locus,
\[
\alpha(n) \;=\; \frac{\Delta_{n}'(u)}{H_{n+1}(u)} - \frac{\Delta_{n-1}'(u)}{H_{n}(u)},
\qquad
\beta(n) \;=\; \frac{H_{n+1}(u)\, H_{n-1}(u)}{H_{n}(u)^{2}},
\]
where $\Delta_{n}'$ is the modified Hankel determinant obtained by
replacing the last row of $H_{n+1}$ with $(m(1),m(2),\ldots,m(n+1))$.
Theorem~\ref{thm:CW} computes the same quantities without forming or
inverting any of these determinants.
\end{theorem}

\begin{proof}
Cramer's rule applied to the orthogonality system yields the displayed
ratios; equality with the recurrence output follows by uniqueness in
Theorem~\ref{thm:CW}.
\end{proof}

\section{Pad\'e numerator and denominator}

\begin{definition}[Reciprocal flip]\label{def:flip}
For $p(x) = \sum_{k=0}^{n} c_{k}\, x^{k}$, let
$p^{\flat}(z) := z^{n} p(1/z) = \sum_{k=0}^{n} c_{k}\, z^{n-k}$.
\end{definition}

\begin{theorem}[Pad\'e denominator from OPS]\label{thm:padedenom}
Let $g(z,u) = \sum_{k\geq 1} a(k,u)\, z^{k}$. The diagonal $[M/M]$
Pad\'e denominator of $g$ is
\[
\hat Q(M,z) \;=\; P(M,\cdot)^{\flat}(z) \;=\; z^{M} P(M, 1/z).
\]
\end{theorem}

\begin{proof}
Diagonal Pad\'e equivalence (Gragg, 1974, Thm.~1) identifies the
denominator with the OPS member of the moment functional under
$x \leftrightarrow 1/z$.
\end{proof}

\begin{theorem}[Pad\'e numerator via second-kind OPS]\label{thm:padenum}
The associated polynomials $P^{(1)}(n,x)$ defined by
$P^{(1)}(-1,x) = -1,\; P^{(1)}(0,x) = 0$ and the same recurrence as
$P(n,x)$ satisfy: $\hat P(M,z) := z^{M} P^{(1)}(M, 1/z)$ is the diagonal
Pad\'e numerator. Equivalently,
\[
P^{(1)}(n,x) \;=\; \Lop_{y}\!\left[ \frac{P(n,x) - P(n,y)}{x - y} \right],
\]
the $\Lop$-divided-difference of $P(n,\cdot)$.
\end{theorem}

\begin{proof}
The divided-difference formula yields a polynomial of degree $n-1$ in
$x$ satisfying the same recurrence with the stated seeds; the Pad\'e
numerator/denominator pair $(\hat P, \hat Q)$ is then verified to
satisfy $g(z) \hat Q(M,z) - z\hat P(M,z) = O(z^{2M+2})$ by orthogonality
of $P(M,\cdot)$ against $1, x, \ldots, x^{M-1}$.
\end{proof}

\begin{definition}[Approximant]\label{def:approx}
$\hat R(M,z,u) := z\, \hat P(M,z)\,/\,\hat Q(M,z)$.
\end{definition}

\begin{theorem}[Pad\'e identity]\label{thm:padeid}
$g(z,u)\, \hat Q(M,z) - z\,\hat P(M,z) = O(z^{2M+2})$ in $\mathbb{C}[u][[z]]$.
\end{theorem}

\begin{proof}
By Theorem~\ref{thm:padedenom}, $\hat Q(M,z) = z^{M} P(M,1/z)$. Then
$g(z) \hat Q(M,z) = \sum_{j} m(j-1) z^{j} \cdot z^{M} P(M,1/z)$, whose
coefficients of $z^{M+1}, \ldots, z^{2M+1}$ are precisely
$\Lop[x^{k} P(M,x)]$ for $k=0,\ldots,M-1$, all zero by orthogonality.
The $z^{1}, \ldots, z^{M}$ coefficients are matched by
$z\hat P(M,z)$ via Theorem~\ref{thm:padenum}.
\end{proof}

\section{Jacobi matrix and spectrum}\label{sec:jacobi}

\begin{definition}[Jacobi matrix]\label{def:Jacobi}
$J(M,u) \in \mathbb{C}[u]^{M\times M}$ is the non-symmetric tridiagonal
matrix with diagonal $(\alpha(0),\ldots,\alpha(M-1))$, super-diagonal
$(1,\ldots,1)$, sub-diagonal $(\beta(1),\ldots,\beta(M-1))$.
\end{definition}

\begin{theorem}[Spectrum = OPS zeros]\label{thm:spectrum}
$\sigma(J(M,u)) = \{x : P(M,x) = 0\}$ as multisets in $\mathbb{C}$.
\end{theorem}

\begin{proof}
$P(M,x) = \det(xI - J(M,u))$ by expansion along the last row of the
companion matrix realization of multiplication-by-$x$ in the basis
$\{P(0,\cdot), \ldots, P(M-1,\cdot)\}$.
\end{proof}

\begin{corollary}[Pole locus of approximant]\label{cor:poles}
The poles of $\hat R(M,\cdot,u)$ are the reciprocals of the zeros of
$P(M,\cdot)$, i.e. the reciprocals of $\sigma(J(M,u))$.
\end{corollary>

\begin{theorem}[Non-symmetric form is forced]\label{thm:nonsymm}
When $\beta(n) \in \mathbb{C}\setminus\mathbb{R}_{>0}$, the
similarity $J \mapsto D J D^{-1}$ with $D = \mathrm{diag}(d_{n})$,
$d_{n}/d_{n-1} = \sqrt{\beta(n)}$, is well-defined only up to branch
choice and produces complex-symmetric, not Hermitian-symmetric, output.
The signed-quasi-definite case admits no real Hermitian Jacobi
representation.
\end{theorem}

\begin{proof}
A Hermitian tridiagonal matrix has real spectrum; $\sigma(J(M,u))$ is
non-real on a Zariski-open set by Lemma~\ref{lem:signed} and
Corollary~\ref{cor:poles}.
\end{proof}

\section{Convergence}

\begin{theorem}[Stahl, 1997]\label{thm:stahl}
Let $\Delta_{g}(u)$ be the Stahl compact of $g(\cdot,u)$. For every
compact $K \subset \mathbb{C} \setminus \Delta_{g}(u)$,
$\hat R(M,z,u) \to g(z,u)$ uniformly on $K$ as $M\to\infty$, with
convergence rate $\limsup_{M} \|g - \hat R(M,\cdot,u)\|_{K}^{1/M}
= \exp(-2/\mathrm{cap}(\Delta_{g}(u),K))$.
\end{theorem}

\begin{theorem}[Real-ray inclusion]\label{thm:realray}
Under the Volterra contraction hypotheses,
$\Delta_{g}(u) \cap \mathbb{R} = \emptyset$, so $[0,\infty) \subset
\mathbb{C}\setminus\Delta_{g}(u)$ and
\[
\hat R(M, t^{\mu}, u) \;\to\; y(t,u) \qquad \text{uniformly on compacts of }[0,\infty).
\]
\end{theorem}

\begin{proof}
The Volterra fixed-point bound $\|y\|_{T} \leq C_{\mu}(T)\,\|P\|$
extends meromorphically to a strip about $\mathbb{R}_{\geq 0}$ in $z$
under the substitution $z = t^{\mu}$; pole-freeness on this strip
forces $\Delta_{g}(u)\cap\mathbb{R} = \emptyset$.
\end{proof}

\begin{corollary}[Spectral condensation]\label{cor:spectralcond}
$\sigma(J(M,u))^{-1} \to \Delta_{g}(u)$ in Hausdorff distance as
$M\to\infty$, where the inversion is taken pointwise.
\end{corollary}

\section{Integrated and fractionally-integrated Riccati}

\begin{theorem}[Integrated Riccati]\label{thm:integ}
$Y(t,u) := \int_{0}^{t} y(s,u)\, ds$ admits
$Y(t,u) = \sum_{k\geq 1} \frac{a(k,u)}{k\mu+1}\, t^{k\mu+1}$. The
reweighted moments $\tilde m(k,u) := m(k,u)/(k\mu+1)$ generate a new
moment functional $\tilde\Lop$ to which Theorem~\ref{thm:CW} applies
verbatim, producing $\tilde S(n,z), \tilde\alpha(n), \tilde\beta(n),
\tilde P(n,x)$ and Pad\'e denominator
$\tilde{\hat Q}(M,z) = \tilde P(M,\cdot)^{\flat}(z)$.
\end{theorem}

\begin{proof}
$\int_{0}^{t} s^{k\mu}\,ds = t^{k\mu+1}/(k\mu+1)$ termwise. The
reweighting is by the Gamma ratio $\Gamma(k\mu+1)/\Gamma(k\mu+2)$,
which is a unit in $\mathbb{C}[u]$ (independent of $u$), preserving
quasi-definiteness.
\end{proof}

\begin{corollary}[Arbitrary fractional integration]\label{cor:fracint}
For $\rho>0$, $I^{\rho} y(t,u)$ is the Pad\'e approximant of the moment
sequence $m(k,u)\, \Gamma(k\mu+1)/\Gamma(k\mu+1+\rho)$. The case
$\rho = 1-\mu$ recovers $I^{1-\mu} y$, occurring in rough Heston.
\end{corollary}

\begin{corollary}[Closure of the OPS family]\label{cor:closure}
The family of moment functionals
$\{\Lop_{\rho} : \rho \in \mathbb{R}_{\geq 0}\}$ defined by
$\Lop_{\rho}[x^{k}] = m(k,u)\, \Gamma(k\mu+1)/\Gamma(k\mu+1+\rho)$ is
closed under the polynomial Chebyshev--Wheeler recurrence: each member
yields its own $(\alpha_{\rho}, \beta_{\rho}, P_{\rho})$ via
Theorem~\ref{thm:CW} without modification.
\end{corollary}

\section{Rough Heston}

\begin{theorem}[El Euch--Rosenbaum, 2019]\label{thm:roughHeston}
The log-price characteristic function in rough Heston is
\[
\varphi_{T}(u) = \exp\!\Bigl(\theta \int_{0}^{T} h(s,u)\,ds + V_{0}\, I^{1-\alpha} h(T,u) \Bigr),
\]
with $h$ solving $D^{\alpha} h = \tfrac{1}{2}(-u^{2}-iu) +
(iu\rho\nu - \lambda) h + \tfrac{1}{2}\nu^{2} h^{2}$, i.e.
$(P,Q,R,\mu) = (\tfrac{1}{2}(-u^{2}-iu),\, iu\rho\nu - \lambda,\,
\tfrac{1}{2}\nu^{2},\, \alpha)$, $\alpha = H + \tfrac{1}{2}$.
\end{theorem}

\begin{algorithm}[Rough Heston via polynomial CW]\label{alg:rH}
Input: $(P,Q,R,\alpha,T,V_{0},\theta,M)$, parameter $u$.
\begin{enumerate}[label=(\arabic*),leftmargin=*]
  \item Build $\{m(k,u)\}_{k=0}^{2M}$ by Lemma~\ref{lem:muntz}.
  \item Build $\{S(n,z)\}_{n=0}^{M}$ in $\mathbb{C}[u][z]$ by
        Theorem~\ref{thm:CW}; record $\alpha(n),\beta(n)$.
  \item Recover $P(M,x), P^{(1)}(M,x)$ from $\alpha,\beta$ via
        Definition~\ref{def:OPS} and Theorem~\ref{thm:padenum}.
  \item Compute $\hat Q(M,z) = P(M,\cdot)^{\flat}(z)$,
        $\hat P(M,z) = P^{(1)}(M,\cdot)^{\flat}(z)$.
  \item Evaluate $h^{[M/M]}(t,u) = t^{\alpha} \hat P(M,t^{\alpha})/\hat Q(M,t^{\alpha})$
        at $t=T$ and on the integration grid.
  \item Apply Theorem~\ref{thm:integ} for $\int_{0}^{T} h$ and
        Corollary~\ref{cor:fracint} with $\rho = 1-\alpha$ for
        $I^{1-\alpha} h(T,u)$.
  \item Return $\exp(\theta \int + V_{0}\, I^{1-\alpha} h(T,u))$.
\end{enumerate}
\end{algorithm}

\begin{theorem}[Algorithm correctness]\label{thm:algcorr}
The output of Algorithm~\ref{alg:rH} converges to $\varphi_{T}(u)$
uniformly on compacts of $u$ avoiding the discriminant locus
$\bigcup_{k\leq 2M}\{H_{k}=0\}$, with rate inherited from
Theorem~\ref{thm:stahl}.
\end{theorem}

\begin{proof}
Steps (1)--(4) are exact in $\mathbb{C}[u]$ by Lemma~\ref{lem:muntz},
Theorem~\ref{thm:CW}, Theorem~\ref{thm:padenum}. Step (5) inherits
Stahl convergence (Theorem~\ref{thm:realray}). Steps (6)--(7) inherit
Theorem~\ref{thm:integ} and Corollary~\ref{cor:fracint}. Composition
of uniformly convergent maps with continuous $\exp$ preserves uniform
convergence on compacts.
\end{proof}

\section{Complexity}

\begin{theorem}[Operation count]\label{thm:complexity}
Algorithm~\ref{alg:rH} performs:
\begin{itemize}[leftmargin=*]
  \item $O(M^{2})$ polynomial multiplications in $\mathbb{C}[u]$ for step (1) (M\"untz convolution);
  \item $O(M^{2})$ polynomial coefficient operations in $\mathbb{C}[u]$ for step (2);
  \item $O(M)$ for step (4) (reciprocal flip is index reversal);
  \item $O(M)$ Horner evaluations for step (5).
\end{itemize}
Total: $O(M^{2})$ in $\mathbb{C}[u]$, equivalent to
$O(M^{2} D)$ scalar operations for $D = \deg_{u}$ working bound.
\end{theorem}

\begin{proof}
Counting per Lemma~\ref{lem:degree} and Corollary~\ref{cor:constructive}.
\end{proof}

\begin{theorem}[Comparison with Hankel route]\label{thm:hankelcost}
The Hankel-determinant computation of
$(\alpha(n),\beta(n))_{n=0}^{M-1}$ via Theorem~\ref{thm:hankelequiv}
costs $\Omega(M^{4})$ scalar operations in $\mathbb{C}[u]$ by direct
LU on $H_{n}$ for each $n\leq M$, and is numerically unstable due to
exponential growth of $\mathrm{cond}(H_{n})$ with $n$. Theorem~\ref{thm:CW}
saves a factor $\Omega(M^{2})$ and avoids the conditioning issue.
\end{theorem}

\begin{proof}
$\mathrm{LU}(H_{n})$ is $O(n^{3})$; summing for $n=1,\ldots,M$ gives
$O(M^{4})$. Conditioning of moment Hankel matrices is classical
(Beckermann, 2000, Thm.~4.2).
\end{proof}

\section{Numerical verification}

\begin{proposition}[Validation]\label{prop:verify}
At working precision $128$ bits and bandwidth $M = 64$:
\begin{enumerate}[label=(\roman*),leftmargin=*]
  \item $|\alpha_{CW}(n,u_{0}) - \alpha_{Hankel}(n,u_{0})| < 10^{-47}$
        and identical for $\beta$, for $n\leq M$ and a grid of $u_{0}\in\mathbb{C}$.
  \item Pad\'e residue:
        $\|[z^{M+1\,:\,2M+1}](g\hat Q - z\hat P)\|_{\infty} < 10^{-53}$.
  \item Free-part match: $[z^{1\,:\,M}](g\hat Q - z\hat P) = 0$ to working precision.
\end{enumerate}
\end{proposition}

\section{Corrigendum}

\begin{proposition}[Lemma A correction]\label{prop:corrig}
$\max_{\mu \in (0,1]} \Gamma(\mu+1)/\Gamma(2\mu+1)$ is attained at
$\mu^{*} = 0.14503474316667\ldots$ with value
$1.03967510771\ldots$, replacing the prior values
$\mu^{*}\approx 0.1339$, $h(\mu^{*})\approx 1.0428$. The $21/20$
uniform bound is unchanged.
\end{proposition}

\begin{proof}
$d/d\mu \log[\Gamma(\mu+1)/\Gamma(2\mu+1)] = \psi(\mu+1) - 2\psi(2\mu+1) = 0$
solved to working precision via Newton on $\psi$.
\end{proof}

\section{Extension to $\mu>1$}

\begin{theorem}[Reduction]\label{thm:munugt1}
For $\mu = n+\nu$, $n=\lfloor\mu\rfloor$, $\nu \in [0,1)$, the
M\"untz basis is $\{t^{n+k\nu}\}_{k\geq 0} \cup \{t^{0},\ldots,t^{n-1}\}$.
The fractional part obeys Lemma~\ref{lem:muntz} with
$\Gamma(k\mu+1)/\Gamma((k+1)\mu+1)$ replaced by
$\Gamma(k\nu+n+1)/\Gamma((k+1)\nu+n+1)$. Theorems~\ref{thm:CW},
\ref{thm:padedenom}, \ref{thm:padenum}, \ref{thm:integ} transfer
verbatim.
\end{theorem}

\begin{proof}
The CW recurrence depends only on the moment sequence, which is
produced by the modified Gamma-ratio recurrence; the proof of
Theorem~\ref{thm:CW} is moment-only and oblivious to origin.
\end{proof}

\begin{remark}[Conjectural status]\label{rem:conj}
Theorem~\ref{thm:realray} for $\mu > 1$ requires reproof of
$\Delta_{g}(u)\cap\mathbb{R} = \emptyset$ via iterated Volterra
contraction; conjectural pending that argument.
\end{remark}

\end{document>
